% TOP_PROBLEM: {"id": 4800014, "title": "Multiple ergodic averages — Problem 14", "category_id": 11, "category": {"id": 11, "name": "dynamical_systems", "display_name": "Dynamical Systems", "description": "Problems about long-term behavior of deterministic systems, Hamiltonian dynamics, and periodic orbits.", "slug": "dynamical-systems", "order_index": 11, "created_at": "2026-07-31T15:26:25.671Z"}, "status": "open", "problem_number": "AMR-047-0014", "set_id": 13}
% FIRST_POSTED: 2026-10-01
% ENTRY_KIND: statement_only
% UnsolvedMath ID 4800014; source catalogue code AMR-047-0014
% !TeX program = lualatex
% Definition and sources only; this file is not an LLM solution attempt.
\documentclass[11pt,a4paper]{article}
\usepackage[margin=25mm]{geometry}
\usepackage{fontspec}
\setmainfont{lmroman10-regular.otf}[BoldFont=lmroman10-bold.otf,
 ItalicFont=lmroman10-italic.otf,BoldItalicFont=lmroman10-bolditalic.otf]
\usepackage{amsmath,amssymb,mathtools}
\usepackage{unicode-math}
\setmathfont{latinmodern-math.otf}
\AtBeginDocument{\let\setminus\smallsetminus}
\usepackage{xurl,hyperref}
\hypersetup{colorlinks=true,urlcolor=blue}
\setcounter{secnumdepth}{0}
\setlength{\parindent}{0pt}
\setlength{\parskip}{5pt}
\setlength{\emergencystretch}{3em}
\begin{document}
\section*{Multiple ergodic averages — Problem 14}\label{4800014}
{\small UnsolvedMath ID 4800014 · AMR-047-0014 · Dynamical Systems · AMR Open Problem Lists}
\subsection{Definitions and mathematical statement}
A generalized polynomial is a real-valued function on
$\mathbb Z$ obtained from the identity function $n\mapsto n$
and real constants by finitely many additions,
multiplications, and applications of the floor function.
For example, $n\mapsto\lfloor\alpha n\lfloor\beta n\rfloor
+\gamma n^2\rfloor$ is such a function. The number and
nesting of operations are finite but otherwise unrestricted.
Call it integer valued if every input integer has an
integer output.

Fix $\ell\ge1$ and integer-valued generalized polynomials
$p_1,\ldots,p_\ell$. Let $T_1,\ldots,T_\ell$ be commuting
invertible measure-preserving transformations of a common
probability space $(X,\mathcal B,\mu)$. Prove that for
all functions $f_i\in L^\infty(\mu)$ the sequence
\[
          \frac1N\sum_{n=1}^N
               \prod_{i=1}^\ell f_i\circ T_i^{p_i(n)}
\]
converges in $L^2(\mu)$ as $N\to\infty$.

Integer outputs may be zero or negative, so negative
iterates are interpreted using invertibility. The
generalized polynomials need not be distinct, have
zero constant term, or satisfy any independence
condition. No ergodicity, mixing, or pointwise
convergence is assumed or requested. The target is
existence of a norm limit for every finite family,
without prescribing its value. Nested floor operations
are essential to the scope: restricting to ordinary
polynomials, or to a single floor applied to a real
polynomial, would omit the general problem.
\subsection{Short English statement}
Do multiple averages along arbitrary integer-valued generalized polynomials always converge in mean for commuting transformations?
\subsection{Sources}
\begin{itemize}
\item[S1] ULAM.AI and the UnsolvedMath Contributors, supplied problems.json, record 4800014 (AMR-047-0014), AMR Open Problem Lists. Catalogue identity and source transcription; CC BY 4.0.
\url{https://huggingface.co/datasets/ulamai/UnsolvedMath}.
\item[S2] Frantzikinakis - Some open problems on multiple ergodic averages (2016), Problem 14. Original source indexed by AMR-047-0014.
\url{https://arxiv.org/abs/1103.3808}.
\end{itemize}
\end{document}
